% Source for the figure of interleave() (sec. 18.2): the Morton codes of
% the points of a grid of 4 x 4 in binary, the bits of x in the even places
% and those of y in the odd ones (left), and the order of the codes on the
% grid (right). Build: make figures (in ../build).
\documentclass[border=4pt]{standalone}
\usepackage[T1]{fontenc}
\usepackage{lmodern}
\usepackage{amsmath,amssymb}
\usepackage{tikz}
\usepackage{lamalunga}
% tikz-ads.tex: the TikZ styles shared by the figures of the memory of the
% manual, \input by each standalone figure: a cell of a vector (cell), the
% cell that an operation touches (touched), a free cell (free), a node of a
% linked structure (node), the labels of indices (index) and of variables
% (var), a pointer (ptr), and \cells and \indices to draw a row of cells.
\usetikzlibrary{arrows.meta,positioning,calc,fit,shapes.multipart}
\definecolor{adsTouched}{RGB}{214,196,236}
\definecolor{adsFree}{RGB}{236,236,236}
\tikzset{
  ads/.style={>={Stealth[length=2mm]},font=\small,
              every node/.append style={inner sep=0pt}},
  cell/.style={draw,minimum width=7mm,minimum height=7mm,
               font=\small\ttfamily,outer sep=0pt},
  touched/.style={cell,fill=adsTouched},
  free/.style={cell,fill=adsFree,text=black!45},
  node/.style={draw,rounded corners=1pt,minimum width=8mm,
               minimum height=6mm,font=\small\ttfamily},
  index/.style={font=\scriptsize\ttfamily,text=black!60},
  var/.style={font=\small\itshape},
  ptr/.style={->,thick}
}
% \cells{name}{x,y}{list}: a row of cells, the first at (x,y), named
% name-0, name-1, ...; each item is text/style, the style one of cell,
% touched and free (e.g., {a/cell,e/touched,/free})
\newcommand{\cells}[3]{%
  \foreach \v/\s [count=\i from 0] in {#3}
    \node[\s] (#1-\i) at ($(#2)+(\i*7mm,0)$) {\v};}
% \indices{name}{n}: the indices 0 .. n-1 under the cells of name
\newcommand{\indices}[2]{%
  \foreach \i in {0,...,\numexpr#2-1\relax}
    \node[index,below=1.2mm of #1-\i] {\i};}

\begin{document}
\begin{tikzpicture}[ads,x=11mm,y=11mm]
  \foreach \x in {0,...,3}
    \foreach \y in {0,...,3} {
      \pgfmathtruncatemacro{\a}{mod(\x,2)}
      \pgfmathtruncatemacro{\b}{mod(\y,2)}
      \pgfmathtruncatemacro{\c}{int(\x/2)}
      \pgfmathtruncatemacro{\d}{int(\y/2)}
      \node[cell,minimum width=11mm,minimum height=11mm]
        at (\x,\y) {\d\c\b\a};
      \pgfmathtruncatemacro{\z}{\a+2*\b+4*\c+8*\d}
      \node[cell,minimum width=11mm,minimum height=11mm] at (\x+5.5,\y) {};
      \coordinate (p\z) at (\x+5.5,\y);
    }
  \draw[edge] (p0) \foreach \z in {1,...,15} { -- (p\z) };
  \foreach \z in {0,...,15} \node[point] at (p\z) {};
  \node[index,anchor=north] at ($(p0)+(0,-0.1)$) {0};
  \node[index,anchor=south] at ($(p15)+(0,0.1)$) {15};
  \foreach \x in {0,...,3} {
    \node[index] at (\x,-0.75) {$x = \x$};
    \node[index] at (\x+5.5,-0.75) {$x = \x$};
  }
  \foreach \y in {0,...,3} {
    \node[index,anchor=east] at (-0.65,\y) {$y = \y$};
    \node[index,anchor=east] at (4.85,\y) {$y = \y$};
  }
\end{tikzpicture}
\end{document}
